\begin{question}
A synchronisation object is required with the following signature
%
\begin{scala}
class SameArgSync{
  def sync(n: Int): Unit = ...
  def shutdown: Unit = ...
}
\end{scala}
%
Each thread~$t$ that requires synchronisation should call~|sync(n)| for
some~|n|.  That thread should synchronise with another thread that calls
|sync| passing in the \emph{same} value for the parameter~|n|.  If there is no
current thread for that value for~|n|, then the thread~$t$ should wait.  If a
subsequent thread~$t'$ calls |sync| with the same value of~|n|, then both $t$
and~$t'$ can return.  

The |shutdown| operation should shut down the synchronisation object, as
usual. 

Implement this class, using the client-server pattern.  Hint: you might use an
instance of |scala.collection.mutable.HashMap| inside the server, mapping
values of~|n| to information about the corresponding thread.

Test your code.  You will need to ensure that your testing system always
terminates. 
\end{question}

%%%%%%%%%%%%%%%%%%%%%%%%%%%%%%%%%%%%%%%%%%%%%%%%%%%%%%%

\begin{answerI}
My code is below.  Each client sends a reply channel, in the normal way.  The
server maintains a map from values of |n| to the corresponding reply channel
for calls of |sync| that have to wait. 
%
\begin{scala}
/** A synchronisation object that allows a thread to synchronise with another
  * thread that submits the same argument to `sync`. */
class SameArgSync{
  /** A reply channel. */
  type ReplyChan = Chan[Unit]

  /** Channel on which clients send messages to the server. */
  private val toServer = new SyncChan[(Int, ReplyChan)]

  /** Synchronise with another thread that submits `n`. */
  def sync(n: Int) = {
    val reply = new OnePlaceBuffChan[Unit]; toServer!((n, reply)); reply?()
  }

  private def server = thread("server"){
    // Map holding pending requests
    val map = new scala.collection.mutable.HashMap[Int, ReplyChan]()
    repeat{
      val (n, reply) = toServer?()
      map.get(n) match{
        case Some(r1) => map -= n; reply!(); r1!()
        case None => map += n -> reply
      }
    }
  }

  fork(server)

  def shutdown = toServer.close()
}
\end{scala}

For testing, we need to ensure that each value of~|n| is used an even number
of times.  As with earlier example, we arrange that each thread in the system
calls |sync| just once, to avoid deadlocks.

In order to test that corresponding calls really are synchronous, we arrange
for each thread to write into a log, before and after the call of |sync|.  At
the end, we traverse the log to check it is correct.  More precisely, the
traversal keeps track of, for each~|n|: (1)~whether there is a blocked call
to~|sync(n)|, and (2)~how many calls to~|sync(n)| may return.
\begin{itemize}
\item When the traversal encounters a call: (a)~if there is no corresponding
  blocked call, then this call is blocked; (b)~if there is a corresponding
  blocked call, then that call is unblocked, and both calls may return. 

\item When the traversal encounters a return, it checks that such a return is
  allowed, and decrements the count of allowed returns.
\end{itemize}
At the end of the traversal, all threads should have returned.  

My code is below, explained by comments.
%
\begin{scala}
object SameArgSyncTest{
  val N = 10  // sync is called on values in the range £[0..N)£. 
  val p = 10 // # workers in each test

  // Events to be stored in the log
  trait LogEvent
  case class CallEvent(n: Int) extends LogEvent
  case class ReturnEvent(n: Int) extends LogEvent

  /** Worker that calls £sync(n)£, logging the calls. */
  def worker(me: Int, sync: SameArgSync, log: Log[LogEvent], n: Int) = 
    thread("worker"+me){
      log.add(me, CallEvent(n)); sync.sync(n); log.add(me, ReturnEvent(n))
    }

  /** Do a single test. */
  def doTest() = {
    val sync = new SameArgSync; val log = new Log[LogEvent](p)
    assert(p%2 == 0); val p2 = p/2   
    val ns = Array.fill(p2)(Random.nextInt(N)) // £worker(i)£ will use £ns(i\%p2)£
    run(|| (for(i <- 0 until p) yield worker(i, sync, log, ns(i%p2))))
    checkLog(log.get); sync.shutdown
  }

  /** Check the contents of the log. */
  def checkLog(log: Array[LogEvent]) = {
    // Bit map indicating those £n£ for which a £sync(n)£ is pending
    val pending = new Array[Boolean](N)
    // Count of how many £sync(n)£s can return
    val canReturn = new Array[Int](N)
    for(i <- 0 until log.length){
      log(i) match{
        case CallEvent(n) => 
          if(pending(n)){ pending(n) = false; canReturn(n) += 2 }
          else pending(n) = true
        case ReturnEvent(n) =>
          assert(canReturn(n) > 0, log.take(i+1).mkString("\n"))
          canReturn(n) -= 1
      }
    }
    assert(pending.forall(_ == false) && canReturn.forall(_ == 0))
  }

  def main(args: Array[String]) = {
    for(i <- 0 until 5000){ doTest(); if(i%100 == 0) print(".") }
    println()
  }
}
\end{scala}
\end{answerI}
